% from ega4-1-fr.tex
\oldpage[IV-1]{226}

Pour abréger, nous appellerons \emph{point maximal} d'un préschéma $X$ tout
point générique d'une composante irréductible de $X$~; si
$X=\operatorname{Spec}(A)$ est affine, les points maximaux de $X$ sont les
\emph{idéaux premiers minimaux} de $A$.

\begin{proposition}[1.1.5]
\label{IV.1.1.5-fr}
Soit $f:X\to Y$ un morphisme quasi-compact. Les conditions suivantes sont
équivalentes~:
\begin{enumerate}
  \item[a)] $f$ est dominant~;
  \item[b)] pour tout point maximal $y$ de $Y$, on a
  $f^{-1}(y)\neq\varnothing$~;
  \item[c)] pour tout point maximal $y$ de $Y$, $f^{-1}(y)$ contient un point
  maximal de $X$.
\end{enumerate}
\end{proposition}

L'équivalence de \emph{a)} et \emph{c)} a été démontrée dans (I, 6.6.5). Il
est clair que \emph{c)} entraîne \emph{b)}~; d'autre part, si $X'$ est une
composante irréductible de $X$ rencontrant $f^{-1}(y)$, $f(X')$ est
irréductible et contient $y$, donc est contenu dans la composante irréductible
$Y'=\overline{\{y\}}$ de $Y$ (0\textsubscript{I}, 2.1.6)~; si $x$ est le
point générique de $X'$, on a donc nécessairement $f(x)=y$
(0\textsubscript{I}, 2.1.5)~; donc \emph{b)} entraîne \emph{c)}.

\begin{proposition}[1.1.6]
\label{IV.1.1.6-fr}
Soient $f':X'\to Y$, $f'':X''\to Y$ deux morphismes, $X$ le préschéma somme
$X'\coprod X''$, $f:X\to Y$ le morphisme égal à $f'$ dans $X'$ et à $f''$
dans $X''$. Pour que $f$ soit quasi-compact, il faut et il suffit que $f'$ et
$f''$ le soient.
\end{proposition}

Cela résulte aussitôt de la définition.

\subsection{Morphismes quasi-séparés}
\label{subsection:IV.1.2-fr}

\begin{definition}[1.2.1]
\label{IV.1.2.1-fr}
Soient $X$, $Y$ deux préschémas. On dit qu'un morphisme $f:X\to Y$ est
\emph{quasi-séparé} (ou que $X$ est \emph{quasi-séparé sur} $Y$) si le
morphisme diagonal
$\Delta_f=(1_X,1_X)_Y:X\to X\times_Y X$ est quasi-compact. On dit qu'un
préschéma $X$ est \emph{quasi-séparé} s'il est quasi-séparé sur
$\operatorname{Spec}(\mathbf Z)$.
\end{definition}

Par définition, tout morphisme séparé est quasi-séparé, $\Delta_f$ étant une
immersion fermée (1.1.2, (i))~; un \emph{schéma} $X$ est un préschéma
quasi-séparé, étant séparé sur $\operatorname{Spec}(\mathbf Z)$ (I, 5.4.1).

\begin{proposition}[1.2.2]
\label{IV.1.2.2-fr}
\begin{enumerate}
  \item[(i)] Tout monomorphisme de préschémas $f:X\to Y$ (en particulier toute
  immersion) est quasi-séparé.
  \item[(ii)] Si $f:X\to Y$ et $g:Y\to Z$ sont deux morphismes quasi-séparés,
  $g\circ f:X\to Z$ est quasi-séparé.
  \item[(iii)] Soient $X$, $Y$ deux $S$-préschémas, $f:X\to Y$ un $S$-morphisme
  quasi-séparé. Alors, pour tout changement de base $g:S'\to S$, le morphisme
  $f_{(S')}:X_{(S')}\to Y_{(S')}$ est quasi-séparé.
  \item[(iv)] Si $f:X\to Y$, $f':X'\to Y'$ sont deux $S$-morphismes quasi-séparés,
  alors
  \[
    f\times_S f':X\times_S X'\to Y\times_S Y'
  \]
  est quasi-séparé.
  \item[(v)] Si $f:X\to Y$, $g:Y\to Z$ sont deux morphismes tels que $g\circ f$
  soit quasi-séparé, alors $f$ est quasi-séparé.
  \item[(vi)] Si $f:X\to Y$ est un morphisme quasi-séparé,
  $f_{\mathrm{red}}:X_{\mathrm{red}}\to Y_{\mathrm{red}}$ est quasi-séparé.
\end{enumerate}
\end{proposition}

\begin{proof}
En vertu de (I, 5.5.12), il suffit de démontrer (i), (ii) et (iii).

L'assertion (i) est triviale, tout monomorphisme étant séparé (I, 5.5.1).
Pour démontrer (iii), on se ramène aussitôt au cas où $Y=S$ (I, 3.3.11) et
l'assertion résulte de ce que
$\Delta_{f_{(S')}}=(\Delta_f)_{(S')}$ (I, 5.3.4) et de (1.1.2, (iii)). Pour
démontrer (ii), considérons les projections $p$ et $q$ de $X\times_Y X$ sur
$X$~; si $i=(p,q)_Z$, on sait que le diagramme
\[
\xymatrix{
  X\times_Y X \ar[r]^-{i} \ar[d]_-{\pi} & X\times_Z X
    \ar[d]^-{f\times_Z f} \\
  Y \ar[r]_-{\Delta_g} & Y\times_Z Y
}
\]
